\lectures{22,23}
\title{Fast Fourier Transform}

\begin{document}

\textbf{Definition (Convolution).} Given sequences $\{a_j\}_{j = 0}^n$ and $\{b_k\}_{k = 0}^m$ of numbers, their \textbf{convolution} $a \ast b$ is:
\[ \{c_\ell\}_{\ell = 0}^{n + m} ~~~~ \text{ given by } ~~~~ c_{\ell} = (a \ast b)_{\ell} := \sum_{j = 0}^{\ell} a_j b_{\ell - j}.\]
This is just polynomial multiplication, or the PMF of the sum of two independent discrete PMFs.

\begin{problem}{FFT}
    Given $\{a_j\}_{j = 0}^m$ and $\{b_k\}_{k = 0}^n$, compute their convolution $a \ast b$ in $O((m + n) \log (m + n))$ time.
\end{problem}

\begin{remark}
    The naive brute-force algorithm runs in $O(mn)$ time.
\end{remark}

Let's black-box a solution to FFT for now.  How might it be useful?

\begin{problem}{String Matching}
    Fix a string $s = (s_0, \dots, s_n) \in \{0, 1\}^{n + 1}$ and a pattern $p = (p_0, \dots, p_m) \in \{0, 1\}^{m + 1}$.  Find the set $H$ of all indices $k$ for which $s_{k + i} = p_i$ for all $i = 0, \dots, m$.
\end{problem}

\begin{solution}{String Matching}
    Switch to a $\pm 1$ encoding, reverse $p$, and compute $s \ast p$.  This is $O(n \log n)$. ~ $\blacksquare$
\end{solution}

Now let's un-black-box the solution to FFT.  How is FFT implemented in $O((m + n) \log (m + n))$ time?

\begin{solution}{FFT}
    It's equivalent to compute $p(x) \cdot q(x)$ for polynomials satisfying $\deg (p) = m$ and $\deg(q) = n$.

    The driving principle behind our solution is the following.
    \begin{quote}
        \textbf{Fact.} Any degree-$n$ polynomial is uniquely determined by its evaluations on any $n + 1$ values.
    \end{quote}
    A good choice of values to evaluate our polynomials on, it turns out, is any set of $(2^t)^{\text{th}}$ roots of unity.
    \begin{quote}
        \textbf{Claim (Discrete Fourier Transform).} Let $n$ be a power of two.  For any polynomial $p$ with degree $(n - 1)$ or lower, we can compute $\{p(\omega^{k}) \}_{k = 0}^n$ in $O(n \log n)$ time.

        \emph{Proof:} Divide and conquer.  Say the runtime is $T(n)$.
        \begin{itemize}
            \item Denote $p_{\text{even}}(x)$ and $p_{\text{odd}}(x)$ so that $p(x) = p_{\text{even}}(x^2) + x \cdot p_{\text{odd}}(x^2)$.
            \item In $2T\left(\frac{n}{2}\right)$ time, compute $\{p_{\text{even}}(\omega^{2k})\}_{k = 0}^{n / 2}$ and $\{p_{\text{odd}}(\omega^{2k})\}_{k = 0}^{n / 2}$.
            \item In $O(n)$ time, we can then compute:
            \[ p(\omega^k) = \begin{cases}
                p_{\text{even}}\left(\omega^{2k}\right) + \omega^k \cdot p_{\text{odd}}\left(\omega^{2k}\right) & \text{if $k < \frac{n}{2}$.} \\
                p_{\text{even}}\left(\omega^{2k - n}\right) - \omega^k \cdot p_{\text{odd}}\left(\omega^{2k - n}\right) & \text{if $k \geq \frac{n}{2}$.}
            \end{cases} \]
        \end{itemize}
        This yields $T(n) = 2T\left(\frac{n}{2}\right) + O(n)$, which solves to $T(n) = O(n \log n)$. ~ $\square$
    \end{quote}
    Thus, we select $\omega = e^{\frac{2\pi i}{T}}$ for some power-of-two $T > m + n$, then compute $\{p(\omega^k) \cdot q(\omega^k)\}_{k = 1}^{T}$ in $O((m + n) \log (m + n))$ time.  It remains to invert the DFT process above to extract $p(x) \cdot q(x)$ from here.
    \begin{quote}
        \textbf{Claim (Inverse Discrete Fourier Transform).} Given the values $\{p(\omega^k)\}_{k = 1}^n$ for some polynomial $p$ with degree $(n - 1)$ or lower, we can compute $p(x)$ in $O(n \log n)$ time.

        \emph{Proof:} The point is that the Inverse DFT is no harder than the DFT.

        Say $p(x) = a_0 + a_1x + a_2x^2 + \dots + a_{n - 1}x^{n - 1}$.  Then the DFT computes a matrix product:
        \[ \begin{bmatrix} p(\omega^0) \\ p(\omega^1) \\ p(\omega^2) \\ \vdots \\ p(\omega^{n - 1}) \end{bmatrix}
        = \underbrace{\begin{bmatrix} 1 & 1 & 1 & \cdots & 1 \\ 1 & \omega & \omega^2 & \cdots & \omega^{n - 1} \\ 1 & \omega^2 & \omega^4 & \cdots & \omega^{2(n - 1)} \\ \vdots & \vdots & \vdots & \ddots & \vdots \\ 1 & \omega^{n - 1} & \omega^{2(n - 1)} & \cdots & \omega^{(n - 1)(n - 1)} \end{bmatrix}}_{W} \begin{bmatrix} a_0 \\ a_1 \\ a_2 \\ \vdots \\ a_{n - 1} \end{bmatrix}. \]
        We effectively wish to find an efficient way to multiply the LHS by $W^{-1}$.  The trick is that $W$ satisfies a really nice property:
        \[ \left(\overline{W} W\right)_{s, t} = \sum_{k = 0}^{n - 1} \overline{W}_{s, k} \cdot W_{k, t} = \sum_{k = 0}^{n - 1} \overline{\omega^{s \cdot k}} \cdot \omega^{k \cdot t} = \sum_{k = 0}^{n - 1} \omega^{k \cdot (t - s)} = \begin{cases} n & \text{ if } s = t. \\
            0 & \text{ otherwise.} \end{cases} \]
        Therefore, $\overline{W}W = nI$, so all we need to do is perform the DFT on the LHS and divide by $n$. ~ $\square$
    \end{quote}
    And that completes the implemntation of FFT! ~ $\blacksquare$
\end{solution}

\begin{remark}
    In 18.701 terms, the DFT matrix $W$ satisfies the property that $\frac{1}{\sqrt{n}}W$ is unitary.
\end{remark}

\end{document}